% =========================================================
% File 3: keyanswers.tex
% =========================================================

\section*{Question 1}

\textbf{Answer: b) one-to-one but not onto.}

Let $f:\mathbb{Z}\to\mathbb{Z}$ be defined by $f(n)=2n+1$.

To prove that $f$ is one-to-one, suppose $f(a)=f(b)$. Then
\[
2a+1=2b+1.
\]
Thus,
\[
2a=2b,
\]
and hence
\[
a=b.
\]
Therefore, $f$ is one-to-one.

However, $f$ is not onto because $f(n)=2n+1$ is always odd. For example, there is no
integer $n$ such that
\[
2n+1=0.
\]
Therefore,
\[
\boxed{\text{$f$ is one-to-one but not onto.}}
\]

\section*{Question 2}

\textbf{Answer: $f$ is one-to-one and onto.}

Let $f:\mathbb{R}\to\mathbb{R}$ be defined by
\[
f(x)=3x-2.
\]

First, we prove that $f$ is one-to-one. Suppose $f(x_1)=f(x_2)$. Then
\[
3x_1-2=3x_2-2.
\]
Thus,
\[
3x_1=3x_2,
\]
so
\[
x_1=x_2.
\]
Therefore, $f$ is one-to-one.

Now, we prove that $f$ is onto. Let $y\in\mathbb{R}$. We need to find
$x\in\mathbb{R}$ such that
\[
f(x)=y.
\]
Solving
\[
3x-2=y
\]
gives
\[
x=\frac{y+2}{3}.
\]
Since $\frac{y+2}{3}\in\mathbb{R}$, such an $x$ exists for every
$y\in\mathbb{R}$.

Therefore, $f$ is onto. Hence,
\[
\boxed{\text{$f$ is one-to-one and onto.}}
\]